\section{Conclusion}

We proposed a progressive depth phased repair algorithm for neural network-based control barrier functions. The key insight is that verification depth determines which violations are detectable, leading to a two-phase strategy that first repairs definitive violations at shallow depth, then progressively increases depth to address uncertain regions.

Our theoretical contributions include:
\begin{itemize}
    \item A formal analysis of the NCBF repair problem as constrained optimization with QP projection
    \item An unbiasedness and variance analysis for the randomized smoothing Jacobian estimator applied to CBF bounds
    \item Convergence criteria based on plateau detection for the progressive depth strategy
\end{itemize}

Experiments on multiple dynamical systems demonstrate consistent improvement over baselines, with the progressive depth strategy providing 5-15\% gains in certified percentage.

\subsection{Limitations and Future Work}

Our method assumes that the dynamics $f(x)$ is Lipschitz continuous and the neural network is differentiable almost everywhere, which holds for standard activation functions but may fail for more exotic architectures. Future work includes:
\begin{enumerate}
    \item Extending the method to handle control-affine systems with $g(x)u$ terms
    \item Investigating adaptive depth scheduling based on violation patterns
    \item Combining with adversarial training for robustness
\end{enumerate}